% Scalar recurrence input for the order-zero RG cutoff estimate.
\section{A scalar decay estimate for cutoff errors}
\label{sec:ex-lattice-flat-rate}

The following elementary lemma converts a one-step contraction with a
polynomially controlled source into a power of the ratio of scales. In an RG
application, the difficult step is to prove that the difference of two
quadratic precisions has a positive coefficient majorant satisfying the
hypotheses below. The lemma here does not prove that RG estimate.

\begin{proposition}[Contraction with a polynomially controlled source]
\label{prop:ex-lattice-flat-rate}
Let $M\in\mathbb N$, and let
\[
 t_0<t_1<\cdots<t_M
\]
be real numbers whose successive gaps satisfy
\[
 \Delta_j:=t_{j+1}-t_j\geq\delta>0
 \qquad(0\leq j<M).
\]
Fix $s\in\mathbb R$, $Z\geq0$, $d\in\mathbb N$, $c_0,\theta>0$, and $A\geq0$.
Put $T=t_M$, choose
\[
 \mu=\min\{c_0,\theta/2\},\qquad
 \lambda=\theta-\mu,\qquad
 R=\max\left\{2,\frac{2(d+1)}{\lambda}\right\},\qquad
 y_j=R+Z+|t_j-s|,
 \]
and set $q=\exp(-\lambda\delta/2)$ and $C=A/(1-q)$. Suppose
$f_0,\ldots,f_M$ and $g_0,\ldots,g_{M-1}$ are nonnegative real numbers such
that, for $0\leq j<M$,
\begin{equation}
 f_j\leq e^{-\theta\Delta_j}f_{j+1}+g_j,
 \qquad e^{c_0(T-t_j)}g_j\leq A y_j^d,
 \label{eq:ex-lattice-flat-rate-recurrence}
\end{equation}
and at the terminal index
\begin{equation}
 f_M\leq A y_M^d.
 \label{eq:ex-lattice-flat-rate-boundary}
\end{equation}
Then, for every $j\in\{0,\ldots,M\}$,
\begin{equation}
 f_j\leq C y_j^d e^{-\mu(T-t_j)}.
 \label{eq:ex-lattice-flat-rate-conclusion}
\end{equation}
If positive scales are defined by $a_j=e^{-t_j}$, the exponential factor in
\eqref{eq:ex-lattice-flat-rate-conclusion} equals $(a_M/a_j)^\mu$.
\end{proposition}

\begin{proof}
The constants in the conclusion depend on $A,d,c_0,\theta,\delta$ but not on
$s$ or $Z$. The choices give $0<\mu\leq c_0$, $\mu\leq\theta/2$, and hence
$\lambda=\theta-\mu\geq\theta/2>0$. Thus $0<q<1$ and $C\geq A$.
Introduce the tilted errors
\[
 F_j=e^{\mu(T-t_j)}f_j,\qquad
 G_j=e^{\mu(T-t_j)}g_j.
\]
Since $t_j\leq T$ and $\mu\leq c_0$, the source bound in
\eqref{eq:ex-lattice-flat-rate-recurrence} implies
\begin{equation}
 0\leq G_j\leq A y_j^d.
 \label{eq:ex-lattice-flat-rate-tilted-source}
\end{equation}
Multiplying the recurrence by $e^{\mu(T-t_j)}$ gives
\begin{equation}
 F_j\leq e^{-\lambda\Delta_j}F_{j+1}+G_j.
 \label{eq:ex-lattice-flat-rate-tilted-recurrence}
\end{equation}

We first control the change in the polynomial weight. The triangle
inequality gives
\[
 y_{j+1}=R+Z+|t_{j+1}-s|\leq y_j+\Delta_j.
\]
Because $y_j\geq R\geq2(d+1)/\lambda$, we have
$d/y_j\leq\lambda/2$. Using $1+x\leq e^x$ for $x\geq0$, we obtain
\begin{equation}
\begin{split}
 e^{-\lambda\Delta_j}y_{j+1}^d
 &\leq y_j^d e^{-\lambda\Delta_j}
       \left(1+\frac{\Delta_j}{y_j}\right)^d\\
 &\leq y_j^d e^{-(\lambda-d/y_j)\Delta_j}\\
 &\leq y_j^d e^{-\lambda\Delta_j/2}
 \leq q y_j^d.
 \end{split}
 \label{eq:ex-lattice-flat-rate-age}
\end{equation}
For $d=0$, the same inequalities hold with the polynomial factor equal to one.

We now argue backward from $M$. At the endpoint,
$F_M=f_M\leq A y_M^d\leq C y_M^d$. Suppose $F_{j+1}\leq C y_{j+1}^d$.
Equations \eqref{eq:ex-lattice-flat-rate-tilted-recurrence},
\eqref{eq:ex-lattice-flat-rate-tilted-source}, and
\eqref{eq:ex-lattice-flat-rate-age} imply
\[
 F_j\leq Cq y_j^d+A y_j^d=(Cq+A)y_j^d=C y_j^d,
\]
where the last equality follows from $C=A/(1-q)$. Induction proves this bound
for every $j$. Undoing the tilt gives
\[
 f_j=e^{-\mu(T-t_j)}F_j
 \leq C y_j^d e^{-\mu(T-t_j)},
\]
as required. Finally,
$e^{-(T-t_j)}=e^{-t_M+t_j}=a_M/a_j$, so the exponential factor is the
stated scale ratio to the power $\mu$.
\end{proof}

\begin{remark}[Scope of the recurrence input]
\label{rem:ex-lattice-flat-rate-scope}
The proposition is a complete scalar argument once the nonnegative recurrence,
its polynomially bounded tilted source, and its endpoint estimate are known.
In the Yang--Mills construction, those data must be obtained from the
scale-by-scale propagation of a neutral curvature representative for the
difference of two quadratic RG precisions. Establishing that propagation
requires the exact Schur-complement, locality, and weighted-kernel hypotheses
of the RG step. This proposition supplies none of them. It therefore proves a
power-rate consequence of that RG input, not the RG input itself or convergence
of the public Wilson moments.
\end{remark}
